% ===========================================================================
%  Planteamiento del proyecto final — ¿En qué minuto se perdió la partida?
%  Héctor Manuel Morán Maldonado
%  Compilar:  latexmk -pdf planteamiento_proyecto_final.tex
% ===========================================================================
\documentclass[11pt,letterpaper]{article}

\usepackage[utf8]{inputenc}
\usepackage[T1]{fontenc}
\usepackage[spanish,es-tabla,es-nodecimaldot]{babel}
\usepackage{lmodern}
\usepackage{microtype}
\usepackage[margin=2.5cm,top=2.6cm,bottom=2.6cm]{geometry}
\usepackage{graphicx}
\usepackage{booktabs}
\usepackage{array}
\usepackage{longtable}
\usepackage{float}
\usepackage{xcolor}
\usepackage{amsmath}
\usepackage{enumitem}
\usepackage{needspace}
\usepackage{caption}
\usepackage{fancyhdr}
\usepackage{titlesec}
\usepackage[hidelinks]{hyperref}

\emergencystretch=2em
\definecolor{azul}{HTML}{2A78D6}
\definecolor{naranja}{HTML}{EB6834}
\definecolor{aqua}{HTML}{1BAF7A}
\definecolor{tinta}{HTML}{1A1A19}
\definecolor{suave}{HTML}{52514E}
\definecolor{papel}{HTML}{F4F6F8}

\captionsetup{font=small,labelfont={bf,color=azul},width=.92\linewidth}
\titleformat{\section}{\Large\bfseries\color{tinta}}{\thesection}{.7em}{}
\titleformat{\subsection}{\large\bfseries\color{suave}}{\thesubsection}{.7em}{}
\titlespacing*{\section}{0pt}{2.2ex plus 1ex minus .2ex}{1.2ex plus .2ex}

\pagestyle{fancy}\fancyhf{}
\renewcommand{\headrulewidth}{.4pt}
\fancyhead[L]{\footnotesize\color{suave}Planteamiento del proyecto final · ¿En qué minuto se perdió la partida?}
\fancyhead[R]{\footnotesize\color{suave}H. M. Morán Maldonado}
\fancyfoot[C]{\footnotesize\color{suave}\thepage}

\newsavebox{\cajaresp}
\newenvironment{respuesta}[1]
  {\par\medskip\setlength{\fboxsep}{10pt}%
   \begin{lrbox}{\cajaresp}%
   \begin{minipage}{\dimexpr\linewidth-2\fboxsep-3pt\relax}%
   \color{tinta}\textcolor{azul}{\textbf{#1}}\par\smallskip\ignorespaces}
  {\end{minipage}%
   \end{lrbox}%
   \noindent\textcolor{azul}{\rule[-\dp\cajaresp]{3pt}{\dimexpr\ht\cajaresp+\dp\cajaresp\relax}}%
   \colorbox{papel}{\usebox{\cajaresp}}\par\medskip}

\newcommand{\dato}[1]{\textbf{\textcolor{azul}{#1}}}
\newcommand{\hecho}{\textcolor{aqua}{\textbf{Hecho}}}
\newcommand{\falta}{\textcolor{naranja}{\textbf{Falta}}}
\newcommand{\parcial}{\textcolor{azul}{\textbf{Parcial}}}

\begin{document}

% ===========================================================================
\begin{titlepage}
\thispagestyle{empty}
\centering
\vspace*{0.6cm}
\includegraphics[width=6.4cm]{assets/logo-up.png}

\vspace{1.1cm}
{\large\color{suave}Especialidad en Ciencia de Datos\par}
\vspace{0.25cm}
{\large\color{suave}Visualización de datos y narración de historias\par}

\vspace{1.5cm}
{\color{azul}\rule{\linewidth}{1.6pt}}
\vspace{0.85cm}

{\Huge\bfseries Planteamiento del proyecto final\par}
\vspace{0.45cm}
{\LARGE \emph{¿En qué minuto se perdió la partida?}\\[.25cm]
 Anatomía de mis derrotas en clasificatoria\par}
\vspace{0.35cm}
{\large\color{suave} Qué se va a presentar, cómo se va a construir y en qué orden\par}

\vspace{0.85cm}
{\color{azul}\rule{\linewidth}{1.6pt}}

\vfill

\begin{minipage}{0.46\linewidth}
  \raggedright
  {\color{suave}\small ALUMNO}\par\smallskip
  {\large Héctor Manuel Morán Maldonado\par}
\end{minipage}\hfill
\begin{minipage}{0.46\linewidth}
  \raggedleft
  {\color{suave}\small PROFESOR}\par\smallskip
  {\large Dr.\ Rodolfo Rubén Álvarez González\par}
\end{minipage}

\vspace{1.1cm}
{\color{suave}\rule{\linewidth}{.4pt}}
\vspace{0.45cm}

{\color{suave}\small
Entregables: presentación oral de máximo 10 minutos (30\,\%) y reporte en PDF
con la construcción de los recursos (40\,\%)\par}
\end{titlepage}

\tableofcontents
\newpage

% ===========================================================================
\section{Resumen y estado actual}

El proyecto reconstruye mis partidas clasificatorias de League of Legends minuto a
minuto, con la línea temporal de la API de Riot, para identificar \textbf{en qué
minuto cambia la dirección de una partida perdida}, si ese minuto forma patrones y si
las derrotas se pueden clasificar por \emph{cómo} ocurrieron.

\begin{respuesta}{Dónde está el proyecto hoy}
Están hechos el planteamiento (presentación de la primera entrega), el EDA piloto y la
política de anonimización sobre 7 partidas públicas de ejemplo, y \textbf{ya está
descargada la muestra propia}: \dato{1\,241 partidas} de cola 420 de mis \dato{seis
cuentas} en LAN. Después del filtrado quedan \dato{1\,195 partidas: 658 victorias y
537 derrotas}. La muestra ya está anonimizada y el barrido de fugas da cero, y los
nueve experimentos (E0--E8) ya corrieron. Están hechos también el reporte final, V1, las
presentaciones y el repositorio público. Este documento conserva el plan tal como se
registró; las cifras y veredictos definitivos son los del reporte final.
\end{respuesta}

\begin{longtable}{p{4.6cm}p{2.0cm}p{8.0cm}}
\caption{Cada punto de la rúbrica, dónde se va a cubrir y en qué estado está.}\\
\toprule
Punto de la rúbrica & Estado & Qué existe o qué falta \\
\midrule
\endfirsthead
\toprule
Punto de la rúbrica & Estado & Qué existe o qué falta \\
\midrule
\endhead
Descripción del problema & \hecho & Láminas 2--8 del planteamiento; se reutilizan. \\
Hipótesis iniciales & \hecho & H1--H3, presentadas en la primera entrega (sección~\ref{sec:hip}). \\
Hipótesis finales & \hecho & Con la muestra propia las tres se mantienen: H1 se precisa, H2 queda no concluyente y H3 se confirma (sección~\ref{sec:prelim}). \\
Bases de datos & \hecho & 1\,241 partidas de seis cuentas descargadas y verificadas (sección~\ref{sec:datos}). \\
Etapa de filtrado & \hecho & Embudo medido: de 1\,241 a 1\,195 partidas (tabla~\ref{tab:embudo}). \\
Preprocesamiento & \hecho & Cuatro tablas, anonimización y medida \(v(t)\) sobre la muestra propia. \\
Modelos / experimentos & \hecho & E0--E8 hechos (sección~\ref{sec:exp}), con sus desviaciones documentadas. \\
Visualizaciones & \hecho & V1--V10 hechas; V1 es una derrota propia elegida por regla (la mediana de A2). \\
Conclusiones & \hecho & En el reporte final; solo afirmaciones de asociación, no causales. \\
Reporte reproducible & \hecho & Reporte final en PDF con el código en anexo, y repositorio público: github.com/hmoran457/lol-cuando-se-pierde-la-partida. \\
\bottomrule
\end{longtable}

% ===========================================================================
\section{Descripción del problema}

Al terminar una partida, el juego muestra «DERROTA» y un marcador: KDA, daño, oro.
Esa es la primera evidencia que todos usamos para explicar el resultado, y es la más
pobre, porque \textbf{describe el final en lugar de explicarlo}. La partida pública
que se usó en el planteamiento lo muestra: el equipo que perdió nunca estuvo detrás
en oro, llegó a +8\,052 en el minuto 22 y todavía iba +1\,034 cuando cayó su Nexo.

\noindent La pregunta del proyecto es:
\begin{quote}
\emph{¿Mis partidas se pierden por un gran error final, o por la acumulación de
decisiones anteriores que nadie registra como error? ¿Y existe un minuto, medible,
en que la derrota deja de ser reversible?}
\end{quote}

\noindent \textbf{Lo que no busca}, y se repetirá en la presentación final: no busca
culpables, no mide estado emocional (\emph{tilt}), no predice partidas futuras y no
demuestra causalidad. El dato es observacional y de equipo: describe las derrotas
\emph{de mis equipos}, no mi aporte individual, que no es separable con esta fuente.
La unidad es \textbf{un jugador en seis cuentas}, todas jugadas por mí en el mismo
servidor (LAN).

% ===========================================================================
\section{Hipótesis}
\label{sec:hip}

\subsection{Hipótesis iniciales (primera entrega)}

El EDA piloto partió de seis hipótesis. En la presentación del planteamiento se
redujeron a \textbf{tres}, a propósito: tres pruebas al 5\,\% dan un 14\,\% de
probabilidad de al menos un falso positivo, y seis daban un 26\,\%. Las tres corren
sobre la línea temporal y forman un solo hilo.

\begin{itemize}[leftmargin=1.4em,itemsep=.45em]
  \item \textbf{H1 · Ventaja desaprovechada.} Con la misma ventaja temprana, mi equipo
        toma menos objetivos por minuto en las partidas que perdemos que en las que
        ganamos.\\
        \emph{Trampa:} «desperdiciamos la ventaja» es una tautología si se mira el
        resultado. Se mide sin mirarlo: objetivos por minuto mientras vamos arriba,
        en derrotas y victorias.
  \item \textbf{H2 · El quiebre coincide con un objetivo.} El minuto en que más cae la
        ventaja coincide con un objetivo mayor del rival (barón, dragón, heraldo,
        larvas, Atakhan) más seguido de lo que explicaría el azar.\\
        \emph{Trampa:} el objetivo marca el momento, no lo causa; y las caídas
        tardías son mayores solo porque hay más oro, así que hay que normalizar.
  \item \textbf{H3 · Punto de no retorno.} Existe un minuto a partir del cual mis
        derrotas ya casi nunca se revierten.\\
        \emph{Trampa:} sesgo de supervivencia. En el minuto 30 solo quedan las
        partidas que llegaron al 30; el denominador se encoge justo donde estaría el
        efecto.
\end{itemize}

\subsection{Hipótesis finales: cómo se van a decidir}

Las hipótesis finales pueden ser las mismas o cambiar, pero el criterio se fija
\textbf{ahora}, antes de ver mis datos, para que un cambio no parezca un ajuste a
conveniencia:

\begin{enumerate}[leftmargin=1.6em,itemsep=.3em]
  \item Una hipótesis \textbf{se mantiene} si la muestra alcanza para contrastarla,
        aunque el resultado sea en contra. Refutar H1 es un resultado, no un fracaso.
  \item Una hipótesis \textbf{se reformula} solo por una razón de dato, documentada:
        por ejemplo, si H3 resulta imposible de ubicar porque después de cierto minuto
        quedan menos de 20 derrotas, se reformula como «¿hasta qué minuto se puede
        siquiera medir un punto de no retorno con esta muestra?». Con la muestra real,
        622 partidas llegan al minuto 30, así que ese caso queda más allá del 30.
  \item \textbf{No se agregan hipótesis} después de ver los datos. Lo que aparezca y
        no estaba previsto se reporta como observación exploratoria, separada de
        las hipótesis.
\end{enumerate}

\noindent En la presentación final, la lámina de hipótesis mostrará las dos versiones
lado a lado (inicial y final), con la razón de cada cambio.

% ===========================================================================
\Needspace{10\baselineskip}
\section{Bases de datos}
\label{sec:datos}

\begin{longtable}{p{3.4cm}p{11.2cm}}
\caption{Fuentes del proyecto.}\\
\toprule
Fuente & Qué aporta \\
\midrule
\endfirsthead
\toprule
Fuente & Qué aporta \\
\midrule
\endhead
Riot Games API, \texttt{match-v5} &
Fuente principal; requiere llave de desarrollador. Por partida: el \emph{registro}
(resultado, campeones, roles, estadísticas finales) y la \emph{línea temporal} (un
cuadro por minuto con el estado de los diez jugadores, más los eventos con marca de
tiempo). Se descargaron las partidas de cola \textbf{420} (clasificatoria
individual) de \textbf{mis seis cuentas}, todas en LAN: 1\,241 pares de JSON,
915\,MB. \\
\addlinespace
Data Dragon 16.17.1 &
CDN estático sin llave. Traduce identificadores a nombres e imágenes (campeones,
objetos). No aporta hechos de partida. \\
\addlinespace
Corpus piloto (7 partidas) &
Ejemplos públicos de \texttt{MingweiSamuel/riotapi-schema}. Se usó para el EDA y la
anonimización; en el proyecto final solo aparece como ejemplo de la pregunta, nunca
como resultado. \\
\bottomrule
\end{longtable}

\begin{respuesta}{La retención se midió, y era más larga de lo supuesto}
El planteamiento suponía que la API conserva la línea temporal \textbf{alrededor de un
año}, según una nota de Riot de 2019. Medido: las 1\,241 líneas temporales pedidas
respondieron, \textbf{ninguna dio 404}, y la más vieja es de \textbf{octubre de 2024},
unos dos años atrás. El supuesto era falso, y en la dirección favorable. El JSON crudo
se guardó de todos modos, porque volver a descargar no daría la misma muestra.
\end{respuesta}

\noindent \textbf{Cómo se descargó:} un script propio
(\texttt{LoL\_Datos/scripts/descargar\_partidas.py}) que pide la lista completa y la de cola 420 explícita,
descarga de la partida más vieja a la más nueva, retoma si la llave caduca o se cae la
red, y respeta los límites de la llave de desarrollo. La llave nunca se escribe a disco.

\noindent \textbf{Seis cuentas, un jugador.} Una sola cuenta daba 41 derrotas, por
debajo del criterio de continuar (\(\sim\)150 partidas). Con las seis cuentas propias
quedan 537 derrotas y el criterio se cumple holgadamente, así que \textbf{se conservan
los cuatro arquetipos}. Las cuentas se identifican como C1--C6, sin nombres.

\begin{table}[H]\centering\small
\caption{La muestra por cuenta, después del filtrado.}
\label{tab:cuentas}
\begin{tabular}{lrrrr}
\toprule
Cuenta & Cola 420 en la lista & Muestra final & Victorias & Derrotas \\
\midrule
C1 & 103 & 102 & 61 & 41 \\
C2 & 552 & 529 & 289 & 240 \\
C3 & 191 & 188 & 99 & 89 \\
C4 & 166 & 156 & 84 & 72 \\
C5 & 141 & 132 & 76 & 56 \\
C6 & 98 & 88 & 49 & 39 \\
\midrule
\textbf{Total} & 1\,251 & \textbf{1\,195} & \textbf{658} & \textbf{537} \\
\bottomrule
\end{tabular}
\end{table}

\noindent La suma de las listas (1\,251) es mayor que las partidas únicas (1\,241)
porque hay partidas que aparecen en dos listas (ver «dúo entre cuentas propias» en el
preprocesamiento).

\noindent \textbf{Cobertura:} de octubre de 2024 a octubre de 2026, \textbf{46 parches}
(14.20 a 16.20), con un hueco casi completo entre febrero y abril de 2026 (13
partidas en tres meses). Mi rol es BOTTOM en 839 de 1\,195 partidas (70\,\%).

% ===========================================================================
\section{Etapa de filtrado}

Sí hubo filtrado, y cada paso se reporta con cuántas partidas elimina, como un
embudo (es la visualización V8, sección~\ref{sec:vis}). Las reglas, y lo que midieron
sobre la muestra:

\begin{enumerate}[leftmargin=1.6em,itemsep=.3em]
  \item \textbf{Lista completa de partidas} de cada cuenta, paginando hasta agotarla.
  \item \textbf{Cola 420 explícita.} Pedir «ranked» a secas mezcla la cola flexible
        (440), que admite grupos de cinco y tiene otra dispersión de nivel.
  \item \textbf{Con línea temporal disponible} (descarta las que dan 404 por
        retención).
  \item \textbf{Diez participantes} y dos equipos. El piloto encontró una partida de
        Arena con 8 jugadores que producía una diferencia de oro falsa \emph{sin
        fallar}.
  \item \textbf{Sin \emph{remakes} ni rendiciones tempranas}, usando
        \texttt{gameEndedInEarlySurrender}, y duración de al menos 5 minutos.
        \textbf{No} se filtra por \texttt{endOfGameResult}: Riot tiene un error
        abierto que marca mal partidas normales, justo en la cola 420.
  \item \textbf{Sin partidas abandonadas:} una rendición antes del minuto 10 casi
        siempre significa que alguien abandonó la partida. No es un \emph{remake} (no
        la marca \texttt{gameEndedInEarlySurrender}) y pasa el filtro de 5 minutos,
        pero la partida no se jugó: se trata como atípica, se excluye y se reporta
        cuántas fueron. La regla se agregó al revisar la primera descarga, antes de
        calcular cualquier resultado.
\end{enumerate}

\begin{table}[H]\centering\small
\caption{Embudo de filtrado medido sobre las seis cuentas.}
\label{tab:embudo}
\begin{tabular}{lrr}
\toprule
Paso & Partidas & Quita \\
\midrule
Lista completa, todas las colas (suma de cuentas) & 2\,442 & — \\
Cola 420, partidas únicas & 1\,241 & — \\
Con registro y línea temporal & 1\,241 & 0 \\
Cola 420 confirmada en el registro & 1\,241 & 0 \\
Diez participantes & 1\,241 & 0 \\
Sin \emph{remake} (\texttt{gameEndedInEarlySurrender}) & 1\,204 & 37 \\
Duración de al menos 5 minutos & 1\,202 & 2 \\
Sin rendición antes del minuto 10 (abandono) & 1\,195 & 7 \\
\midrule
\textbf{Muestra final} & \textbf{1\,195} & \textbf{46} \\
\bottomrule
\end{tabular}
\end{table}

\noindent Las 1\,241 partidas tienen \texttt{endOfGameResult = GameComplete}, así que
no filtrar por ese campo no cambió nada esta vez; la regla se mantiene por el error
de Riot que afecta a la cola 420.

\noindent \textbf{Lo que no se filtra:} las victorias. Entran como grupo de control,
porque toda partida tiene su mayor caída de oro, también las ganadas. Sin ellas, el
análisis describiría cómo son las partidas de League, no cómo se pierden las mías.

% ===========================================================================
\section{Preprocesamiento}

\begin{enumerate}[leftmargin=1.6em,itemsep=.35em]
  \item \textbf{Anonimización primero.} Antes de que cualquier dato llegue a un CSV,
        una gráfica o el repositorio, se aplica la política de la actividad de
        anonimización: nombres suprimidos, \texttt{puuid} y \texttt{matchId}
        seudonimizados con HMAC, fecha generalizada a mes. Medido sobre la muestra
        propia: el barrido busca 9\,852 \texttt{puuid}, 10\,473 nombres y 1\,241
        \texttt{matchId}, los encuentra todos en el crudo (control negativo) y
        \dato{ninguno} en la salida. Con 1\,241 partidas, quitar el campeón deja
        solo el 1\,\% de participantes únicos (eran 63\,\% en el piloto, como
        advertía el reporte de anonimización); la composición de diez campeones
        sigue siendo única en todas, por eso al repositorio solo van tablas a nivel
        equipo.
  \item \textbf{Cuatro tablas} desde el JSON anidado: \texttt{partidas},
        \texttt{participantes}, \texttt{cuadros} (jugador × minuto) y
        \texttt{eventos}. Sobre la muestra propia: 1\,241 partidas, 12\,410
        participantes y más de un millón de eventos.
  \item \textbf{Cuál de los diez soy yo.} Las seis cuentas se marcan antes de
        seudonimizar (\texttt{es\_propio} y \texttt{cuenta} = C1--C6), porque después
        del HMAC ya no se puede saber.
  \item \textbf{Dúo entre cuentas propias.} En \dato{13 partidas} aparecen dos de mis
        cuentas, siempre en el \emph{mismo} equipo: otra persona jugaba una de ellas en
        dúo conmigo. Regla: \textbf{yo soy la cuenta en BOTTOM} (en las 13 hay
        exactamente una), y la otra cuenta es un compañero más. Esas 13 partidas
        tampoco son independientes del resto, y se reporta.
  \item \textbf{Equipo de cada jugador} cruzado desde el registro: la línea temporal
        solo trae \texttt{participantId} y \texttt{puuid}.
  \item \textbf{Diferencia de oro} calculada, porque no es un campo de la API: suma de
        \texttt{totalGold} de un equipo menos el otro, cuadro por cuadro.
  \item \textbf{Medida relativa}
        \(v(t) = \dfrac{\text{diferencia de oro}}{\text{oro total de los diez}}\).
        2\,467 de oro son el 7.7\,\% del mapa en el minuto 10 y el 2.3\,\% en el 28;
        un umbral fijo en oro diría que son lo mismo.
  \item \textbf{Cuadros indexados por posición}, nunca por
        \texttt{timestamp // 60000}: las marcas se desvían hasta +582\,ms y esa
        división junta dos minutos en la misma casilla.
  \item \textbf{Estructuras atribuidas por \texttt{killerId}}, no por
        \texttt{teamId}, que en \texttt{BUILDING\_KILL} nombra al equipo que
        \emph{perdió} la torre. Verificado sobre la muestra propia: las torres por
        equipo cuadran contra el conteo oficial del registro en \dato{1\,195 de
        1\,195} partidas. Las que derriban los súbditos (\texttt{killerId = 0}) se
        atribuyen al equipo contrario al \texttt{teamId}.
  \item \textbf{Vocabulario de eventos desde mis partidas}, no desde la
        documentación. Medido: 20 tipos de evento, y entre los objetivos mayores
        aparecen las larvas del vacío (\texttt{HORDE}, 4\,725) y Atakhan (647), que la
        documentación no lista. Ambos entran en la lista de objetivos de H2 y del
        arquetipo A3.
  \item \textbf{Campos opcionales con acceso defensivo}, con tasa de ausencia
        reportada.
\end{enumerate}

% ===========================================================================
\section{Modelos y experimentos}
\label{sec:exp}

El encuadre no es «seis pruebas de hipótesis», sino una \textbf{taxonomía descriptiva
pre-registrada, con intervalos}: reglas escritas antes de ver la distribución, y cada
resultado con su incertidumbre. Las proporciones se reportan con intervalos de
\textbf{Wilson}, no de Wald, que con muestras pequeñas se equivoca hasta por un factor
de dos en el límite inferior.

\begin{longtable}{p{1.0cm}p{4.2cm}p{9.4cm}}
\caption{Experimentos, en el orden en que se corren.}\\
\toprule
Id & Experimento & Método \\
\midrule
\endfirsthead
\toprule
Id & Experimento & Método \\
\midrule
\endhead
E0 & Tamaño de muestra & \textbf{Hecho.} Embudo de la tabla~\ref{tab:embudo}: 1\,195 partidas, 537 derrotas. Se conservan los cuatro arquetipos. \\
\addlinespace
E1 & Definir el minuto de quiebre & \textbf{Hecho.} D1 (mayor caída de \(v\) en un cuadro) contra D2 (último cruce de \(v\) hacia abajo). Difieren por más de 2 minutos en el \dato{67\,\%} de las derrotas (IC 95\,\%: 63--71). D1 resultó inservible: cae en el minuto 5 de mediana \emph{igual en victorias que en derrotas}, porque al inicio hay poco oro y una sola muerte mueve mucho \(v\). \textbf{D2 pasa a ser la definición principal} (ver desviaciones). \\
\addlinespace
E2 & H1 · objetivos con ventaja & \textbf{Hecho.} Objetivos propios por minuto mientras hay ventaja real (\(v \ge 5\)\,\% y 1\,500 de oro), con bootstrap por \emph{partida}. Agregado: 0.54 en derrotas contra 0.87 en victorias. Por fase, para comparar la misma ventaja: antes del minuto 15, 0.47 contra 0.54 (intervalos traslapados); entre el 15 y el 25, \dato{0.61 contra 0.95}; después del 25, \dato{0.58 contra 1.10}. \\
\addlinespace
E3 & H2 · quiebre y objetivo & \textbf{Hecho.} El 39.6\,\% de los quiebres (D2) de derrotas cae a \(\pm 1\) minuto de un monstruo épico del rival. Contra la nula registrada (45.6\,\% esperado) quedaría \emph{por debajo} del azar, pero esa nula pone objetivos antes del minuto 5, donde no existen. Contra la nula corregida (35.9\,\% esperado) queda por encima: \dato{p = 0.017}. En las victorias la coincidencia es mayor todavía (28\,\% contra 19\,\%). \\
\addlinespace
E4 & H3 · punto de no retorno & \textbf{Hecho.} De las partidas abajo con \(v \le -5\)\,\%, terminó ganando el 37\,\% en el minuto 5, el 24\,\% en el 10, el 20\,\% en el 15, el 12\,\% en el 20, el 8\,\% en el 25 y el \dato{4\,\%} en el 30 (IC 95\,\%: 2--8, 154 en riesgo). La caída es gradual; desde el minuto 27 remonta menos del 5\,\% en todos los minutos siguientes. \\
\addlinespace
E5 & Arquetipos A1--A4 & \textbf{Hecho.} Reglas registradas, en orden A4 → A2 → A1 → A3, sin ajustar umbrales. De 537 derrotas: A1 \dato{19\,\%} (16--23), A4 \dato{15\,\%} (12--18), A2 \dato{10\,\%} (8--13), A3 \dato{4\,\%} (3--6) y \dato{sin clasificar 52\,\%} (48--56). El residuo es la barra más alta: es el criterio de refutación de V6. \\
\addlinespace
E6 & Sensibilidad & \textbf{Hecho.} Con el umbral en 4\,\% cambian de arquetipo 64 derrotas (12\,\%); con 6\,\%, 83 (15\,\%). A3 es el más inestable: 10, 23 y 49 derrotas con 4, 5 y 6\,\%. \\
\addlinespace
E7 & Estabilidad & \textbf{Hecho.} Muestra partida por mes (hasta mayo de 2025: 209 derrotas; después: 328). Mediana de D2: minuto 12 contra 14, no se mueve más de 2 minutos. Reparto de arquetipos igual en ambas mitades (\(\chi^2\), p = 0.71). \\
\addlinespace
E8 & Cuentas y rol & \textbf{Hecho.} Las seis cuentas no se reparten distinto entre arquetipos (permutación, \dato{p = 0.19}), así que se juntan. Solo con BOTTOM (839 partidas, 366 derrotas) todo se repite: arquetipos con diferencias de 3 puntos o menos, H1 por fase 0.42/0.60/0.60 contra 0.55/0.96/1.12, H2 p = 0.034, H3 4\,\% en el minuto 30. Mezclar roles no cambia las conclusiones. \\
\bottomrule
\end{longtable}

\subsection{Desviaciones del plan, documentadas}

Cada cambio respecto a lo escrito antes de ver los datos se declara aquí, con su razón.
Tres de ellos (la definición de E1, la nula de E3 y la comparación por fase de H1) se
decidieron \emph{después} de ver los resultados, al verificar supuestos, y uno ---la nula
de E3--- cambia la conclusión a favor de H2. Por eso, en esos tres casos se reportan
siempre las dos versiones, la registrada y la corregida, para que el lector juzgue.

\begin{enumerate}[leftmargin=1.6em,itemsep=.3em]
  \item \textbf{D2 sustituye a D1 como definición principal del quiebre.} Razón de
        dato (E1): D1 no distingue derrotas de victorias y se concentra en los primeros
        5 minutos. D2 ya estaba registrada como control; no es una regla nueva. D1 se
        reporta como la definición que falló.
  \item \textbf{La nula de E3 se corrige.} La registrada desplazaba los objetivos por
        toda la partida; ningún monstruo épico existe antes del minuto 5 (el primero de
        la muestra cae en el 5.4) y 163 de 535 quiebres de derrotas ocurren antes. Se
        reportan las dos nulas lado a lado (V4).
  \item \textbf{H1 se compara por fase de la partida.} La hipótesis dice «con la misma
        ventaja», y la tasa agregada no lo cumple: al ganar, la ventaja es mayor y más
        tardía. La definición de la tasa no cambia; se reporta también la agregada.
  \item \textbf{A3 usa D2 en lugar de «la mayor caída de \(v\)».} Consecuencia directa
        de E1, decidida antes de correr E5. Las demás reglas de los arquetipos siguen
        la implementación de referencia escrita antes de ver los datos
        (\texttt{probe\_umbrales.py}), incluido el piso de 1\,500 de oro en A1; la
        lista de objetivos de A3 sigue la lámina (con larvas y Atakhan).
  \item \textbf{E8 añade un valor p por permutación}, porque la tabla de cuentas tiene
        celdas con frecuencia esperada menor a 5 (mínima de 1.7) y la aproximación
        \(\chi^2\) no es confiable ahí. Las dos dan lo mismo (0.19).
  \item \textbf{Precisiones de construcción}, fijadas antes de correr cada experimento:
        las placas de torre no cuentan como objetivo; cada larva es un evento; los 86
        monstruos que no tomó ningún equipo se excluyen; un evento se juzga con \(v\) del
        cuadro de su minuto.
\end{enumerate}

\subsection{Propuesta de hipótesis finales}
\label{sec:prelim}

Aplicando los criterios de la sección~\ref{sec:hip} (se mantiene lo que la muestra
alcanza a contrastar, se reformula solo por razón de dato, no se agrega nada):

\begin{longtable}{p{1.0cm}p{6.4cm}p{7.2cm}}
\toprule
 & Versión final & Qué pasó y por qué \\
\midrule
\endfirsthead
\toprule
 & Versión final & Qué pasó y por qué \\
\midrule
\endhead
H1 & \emph{A partir del minuto 15}, con ventaja real, mi equipo toma menos objetivos por minuto en las derrotas que en las victorias. & \textbf{Se mantiene y se precisa.} Antes del 15 la diferencia no se distingue del ruido (0.47 contra 0.54); desde el 15, sí (0.61 contra 0.95). Es asociación: tomar objetivos también produce oro. \\
\addlinespace
H2 & El quiebre (D2: el último minuto en que mi equipo pasa a ir abajo en oro) coincide con un monstruo del rival más de lo que explicaría el azar. & \textbf{Se mantiene; no concluyente.} Con la definición registrada (D1) no se distingue del azar (26.4\,\% contra 25.0\,\%, p = 0.18); con D2 solo supera a la nula corregida (39.6\,\% contra 35.9\,\%, p = 0.017, en el límite de \(\alpha = 0.05/3\)). En las victorias la coincidencia es mayor: describe el ritmo de cualquier partida, no cómo se pierde. \\
\addlinespace
H3 & Existe un minuto a partir del cual mis derrotas ya casi nunca se revierten. & \textbf{Se mantiene y se confirma.} Con «casi nunca» como menos del 5\,\% sostenido (umbral fijado después de ver los datos y declarado como desviación), es el minuto 27 (4.0\,\%, IC 95\,\%: 2.1--7.8, 198 en riesgo). Antes la probabilidad baja de forma gradual (37\,\% en el 5, 20\,\% en el 15), y por eso el minuto depende del umbral (22 con 10\,\%). \\
\bottomrule
\end{longtable}

\noindent \textbf{Los arquetipos describen la mitad.} A1, A4 y A2 juntas explican el
44\,\% de las derrotas; A3 casi no existe (4\,\%) y es inestable. El 52\,\% restante
no cae en ninguna regla, y no se ajustaron umbrales para esconderlo.

\noindent \textbf{Observación exploratoria, no hipótesis:} las derrotas sin clasificar
son \emph{partidas parejas que se deciden tarde}. En la mediana, el 78\,\% de la
partida transcurre con \(|v| < 5\)\,\%, el 75\,\% nunca baja de \(-10\)\,\%, su
quiebre llega en el minuto 19 y duran más (33 minutos contra 30 de las clasificadas).
Se reporta como patrón a explorar, no como un quinto arquetipo.

\subsection{Aritmética de la muestra}

\noindent \textbf{La aritmética, ahora con la muestra real:} con 537 derrotas, un
arquetipo que agrupe el 30\,\% trae un intervalo de Wilson de [0.26, 0.34], y uno del
20\,\%, [0.17, 0.24]. Dos arquetipos al 30\,\% y al 20\,\% difieren 0.10 \(\pm\) 0.05,
así que diferencias de ese tamaño ya se pueden ordenar, lo que con una sola cuenta (41
derrotas) no era posible. \textbf{Con los resultados reales}, el orden se sostiene solo
en parte: A3 es claramente el más raro, pero los intervalos de A1 (16--23\,\%) y A4
(12--18\,\%) se traslapan, así que no se afirma cuál de los dos es más común. Cada
cuenta por separado sigue siendo débil: C6 tiene 39 derrotas.

% ===========================================================================
\section{Visualizaciones}
\label{sec:vis}

Cada visualización existe para soportar o refutar algo concreto. Si no hay hipótesis
detrás, no entra a la presentación (puede ir al reporte como exploración).

\begin{longtable}{p{0.8cm}p{4.0cm}p{2.0cm}p{7.8cm}}
\caption{Visualizaciones planeadas.}\\
\toprule
Id & Visualización & Sirve a & Qué la haría refutar \\
\midrule
\endfirsthead
\toprule
Id & Visualización & Sirve a & Qué la haría refutar \\
\midrule
\endhead
V1 & Curva de oro de una partida, con cortina animada y objetivos marcados & Problema & — (es el gancho). Se sustituye la partida pública por una mía: la \emph{mediana} de su arquetipo, elegida al final. \\
\addlinespace
V2 & \(v(t)\) de todas las derrotas, en paneles pequeños por arquetipo, con la mediana resaltada & E5 & Que los paneles se vean iguales entre sí: los arquetipos no separarían nada. \\
\addlinespace
V3 & Objetivos por minuto con ventaja: punto e intervalo, derrotas contra victorias & H1 & Intervalos que se traslapan por completo. \\
\addlinespace
V4 & Distribución nula de la permutación, con la proporción observada como línea & H2 & Que la línea caiga dentro del cuerpo de la distribución nula. \\
\addlinespace
V5 & Probabilidad de remontar por minuto, banda de Wilson, y barras de partidas en riesgo debajo & H3 & Que la curva no baje, o que la banda sea tan ancha que no permita ubicar ningún minuto. \\
\addlinespace
V6 & Frecuencia de arquetipos con intervalo de Wilson, incluido «sin clasificar» & E5 & Que «sin clasificar» sea la barra más alta. \\
\addlinespace
V7 & Mapa de calor umbral × arquetipo (4\,\%, 5\,\%, 6\,\%) & E6 & Que las partidas salten de arquetipo con un punto de umbral. \\
\addlinespace
V8 & Embudo de filtrado, con cuántas partidas quita cada paso & Datos & — (transparencia). \\
\addlinespace
V9 & Desacuerdo entre las dos definiciones de quiebre (dispersión, minuto contra minuto) & E1 & Puntos lejos de la diagonal: «el minuto de quiebre» dependería de la definición. \\
\addlinespace
V10 & Frecuencia de arquetipos por cuenta (C1--C6), con intervalo de Wilson & E8 & Que una cuenta se separe del resto: entonces no se pueden juntar sin decirlo. \\
\bottomrule
\end{longtable}

\noindent Se mantiene la convención visual del planteamiento: \textcolor{azul}{azul}
es el equipo propio y \textcolor{naranja}{naranja} el rival; lo medido lleva etiqueta
azul y lo ilustrativo, etiqueta ámbar.

% ===========================================================================
\section{Conclusiones: qué se va a poder afirmar}

Las conclusiones aún no existen, pero su \emph{forma} se fija ahora:

\begin{itemize}[leftmargin=1.4em,itemsep=.3em]
  \item Se dirá «coincide con», nunca «causa». El barón no \emph{causó} el desplome:
        la pelea ocurre antes que el objetivo.
  \item Cada porcentaje irá con su intervalo y su \(n\).
  \item Se reportarán también los patrones que \textbf{no} se resuelven con esta
        muestra; que H3 resulte imposible de ubicar sería un resultado legítimo.
  \item Las conclusiones son sobre \emph{un jugador} (seis cuentas, sobre todo en
        BOTTOM) en \emph{una ventana de 46 parches}: no se generalizan a otros
        jugadores, roles ni niveles.
\end{itemize}

% ===========================================================================
\section{Plan de la presentación (10 minutos)}

\begin{longtable}{p{0.6cm}p{5.0cm}r p{7.0cm}}
\caption{Guion de tiempos. Total: 10 minutos.}\\
\toprule
\# & Bloque & Min. & Idea que hay que dejar \\
\midrule
\endfirsthead
\toprule
\# & Bloque & Min. & Idea que hay que dejar \\
\midrule
\endhead
1 & Problema: «DERROTA» y la curva (V1) & 1.5 & El marcador describe el final; la curva muestra el quiebre. \\
2 & Hipótesis iniciales y finales & 1.0 & Las tres, lado a lado con su versión final y la razón de cada cambio. \\
3 & Bases de datos y filtrado (V8) & 1.0 & De dónde sale el dato y cuántas partidas sobreviven a cada filtro. \\
4 & Preprocesamiento & 1.0 & La diferencia de oro no existe: se construye. Y se anonimiza antes de todo. \\
5 & Experimentos & 1.0 & Taxonomía pre-registrada, permutación y Wilson; sin ajustar umbrales. \\
6 & H1 (V3) & 1.0 & Resultado, a favor o en contra. \\
7 & H2 (V4) & 1.0 & ¿Más que el azar? \\
8 & H3 (V5) & 1.0 & Dónde está el punto de no retorno, o hasta dónde se puede medir. \\
9 & Arquetipos (V2, V6) & 1.0 & Las formas de perder y cuántas quedan sin clasificar. \\
10 & Conclusiones & 0.5 & Qué se puede afirmar y qué no. \\
\midrule
 & \textbf{Total} & \textbf{10.0} & \\
\bottomrule
\end{longtable}

\noindent Se reutiliza el motor de la presentación del planteamiento (HTML sin
conexión, cortina animada, riel de secciones). El riel se cambia a las secciones de
esta rúbrica.

% ===========================================================================
\section{Plan del reporte y reproducibilidad}

El reporte cubre los mismos puntos con más profundidad, más todo lo que no cabe en 10
minutos: los dos criterios de quiebre completos, las tablas de sensibilidad, el
embudo con cifras y las observaciones exploratorias. Para que alguien con
conocimientos de análisis de datos pueda reproducirlo:

\begin{itemize}[leftmargin=1.4em,itemsep=.3em]
  \item \textbf{Repositorio público} con los scripts de descarga, anonimización,
        análisis y figuras; un \texttt{README} con el orden exacto de ejecución; y
        las versiones de las bibliotecas fijadas en \texttt{requirements.txt}.
  \item \textbf{Datos en el repositorio:} solo las tablas anonimizadas a nivel equipo
        (\texttt{partidas}, \texttt{oro\_equipo}, \texttt{eventos\_equipo}). El JSON
        crudo y la tabla por participante \textbf{no} se suben: la actividad de
        anonimización midió que esta última es única fila por fila. Quien tenga su
        propia llave puede regenerar todo con el script de descarga.
  \item \textbf{\texttt{.gitignore} antes del primer \texttt{git init}.} Debe excluir
        los JSON crudos desde el inicio: una vez que entran a un commit, quedan en el
        historial aunque después se borren.
  \item \textbf{Semillas fijas} en el bootstrap y la permutación, reportadas.
  \item \textbf{Todas las cifras del texto salen de un archivo de resultados} generado
        por los scripts (como ya hacen el EDA y la anonimización), para que el texto
        no se desincronice de los datos.
  \item \textbf{Código de internet referenciado} en el lugar donde se usa y en la
        sección de referencias.
  \item \textbf{Anexo} con el código de las visualizaciones, o la liga al
        repositorio público.
\end{itemize}

% ===========================================================================
\section{Cronograma}

\begin{longtable}{p{2.0cm}p{12.6cm}}
\caption{Orden de trabajo. Las semanas 1 a 4 ya están hechas.}\\
\toprule
Semana & Trabajo \\
\midrule
\endfirsthead
\toprule
Semana & Trabajo \\
\midrule
\endhead
1 & \hecho: descarga de seis cuentas (E0), embudo de filtrado, anonimización de la muestra propia. Se continúa con cuatro arquetipos. \\
2 & \hecho: medida \(v(t)\) y E1 (D2 pasa a ser la definición principal); V8 y V9. \\
3 & \hecho: E2--E4 (H1, H2, H3); sus cifras observadas coinciden con un cálculo independiente hecho directamente desde el JSON crudo; V3--V5. \\
4 & \hecho: E5--E8 (arquetipos, sensibilidad, estabilidad, cuentas y rol); V2, V6, V7, V10; propuesta de hipótesis finales. \\
5 & \parcial: V1 (mediana de A2) y reporte final en PDF hechos; falta el repositorio público. \\
6 & Presentación de 10 minutos y ensayo cronometrado. \\
\bottomrule
\end{longtable}

% ===========================================================================
\section{Riesgos}

\begin{longtable}{p{5.0cm}p{9.6cm}}
\toprule
Riesgo & Respuesta \\
\midrule
\endfirsthead
\toprule
Riesgo & Respuesta \\
\midrule
\endhead
\textcolor{suave}{Muestra insuficiente} & \textcolor{suave}{Resuelto: con seis cuentas hay 537 derrotas.} \\
\addlinespace
\textcolor{suave}{La retención borra partidas} & \textcolor{suave}{Resuelto: retención medida de al menos dos años, y el crudo ya está en disco.} \\
\addlinespace
\textcolor{suave}{Las cuentas no son comparables} & \textcolor{suave}{Descartado: E8 no encuentra diferencias entre cuentas (p = 0.19).} \\
\addlinespace
\textcolor{suave}{Mezcla de roles (BOTTOM en el 70\,\%)} & \textcolor{suave}{Descartado: solo con BOTTOM los resultados se repiten (E8).} \\
\addlinespace
Residuo «sin clasificar» grande & \textbf{Ocurrió: 52\,\%.} Se reporta como resultado, sin ajustar umbrales, y se describe como observación exploratoria. \\
\addlinespace
A3 inestable ante el umbral & Se reporta E6 junto a A3; no se afirma nada sobre su frecuencia más allá de «rara». \\
\addlinespace
Efecto de parche (46 parches, con un hueco de febrero a abril de 2026) & Se reconoce, no se controla; se reporta cuántas partidas hay por parche. \\
\addlinespace
Partidas no independientes (rachas, sesiones, 13 dúos entre cuentas propias) & Bootstrap por partida; las conclusiones se plantean como descriptivas. \\
\bottomrule
\end{longtable}

\vspace{.5cm}
\noindent{\small\color{suave} Este proyecto no está avalado por Riot Games y no
refleja las opiniones de Riot Games ni de nadie involucrado oficialmente en la
producción o administración de League of Legends. League of Legends y Riot Games
son marcas registradas de Riot Games,~Inc.}

\end{document}
